\chapter{Теоретический раздел}

\section{Постановка задачи}

Целью работы является разработка программной модели туристического агентства,
которая поддерживает ведение каталога туров, регистрацию клиентов,
оформление бронирований, расчет скидок и формирование маркетинговых
предложений. При проектировании системы важно обеспечить расширяемость
бизнес-логики, явное описание жизненного цикла брони и безопасную обработку
ошибок.

В качестве базовых сущностей предметной области используются тур, клиент и
бронирование. Для каждой сущности определяются ограничения корректности
данных. Такой подход позволяет переносить основные инварианты предметной
области непосредственно в модель и упрощает дальнейшую реализацию сервисного
слоя.

\section{Примененные паттерны проектирования}

При реализации системы используются как поведенческие, так и архитектурные
паттерны. Их применение позволяет изолировать изменяемые правила и снизить
связность между модулями~\cite{gamma1995,fowler2002}.

\subsection{Strategy}

Паттерн Strategy используется для расчета скидок. Логика определения размера
скидки вынесена в отдельные политики, например скидку по уровню лояльности и
групповую скидку. Итоговый расчет выполняется объектом
\texttt{PriceCalculator}, который работает через абстракцию политики, а не
через жестко заданные условия. Такое решение позволяет добавлять новые правила
ценообразования без изменения кода сервиса бронирования.

\subsection{State}

Жизненный цикл бронирования описывается паттерном State. Для брони
поддерживаются состояния <<Новая>>, <<В работе>>, <<Готова>>, <<Выдана>> и
<<Отменена>>. Каждое состояние инкапсулирует допустимые переходы и запрещает
некорректные операции. Благодаря этому правила переходов не размазываются по
сервисам и модели, а сосредотачиваются в одном месте.

\subsection{Observer}

Паттерн Observer применяется для обработки уведомлений. Издатель событий
уведомляет подписчиков о смене статуса брони, поступлении оплаты и истечении
срока действия брони. Один наблюдатель сохраняет уведомления в журнале, другой
собирает статистику. Такое разделение позволяет добавлять новые реакции на
события без модификации логики бронирования.

\subsection{Unit of Work и Repository}

Для безопасного выполнения сложных операций используется паттерн Unit of Work.
При создании или отмене бронирования система регистрирует компенсирующие
действия и при ошибке возвращает данные в согласованное состояние. Хранение
данных организовано через репозитории, которые изолируют прикладную логику от
конкретного механизма доступа к данным~\cite{fowler2002}. В учебной версии
используются реализации in-memory, что упрощает тестирование и демонстрацию
поведения.

\section{Архитектурные выводы}

Выбранный набор паттернов хорошо соответствует поставленной задаче. Поведение,
которое меняется независимо, выделено в отдельные классы. Это упрощает
расширение системы, делает тесты более локальными и уменьшает риск
регрессионных ошибок при изменении бизнес-правил.
